% Reviewer xYk4 (W2), small-scale version. Code: text_to_image/_launch_bit_small_scale.sh.
% Fill \tbd{} from $PREFIX_DIR/BiB_results/bit_small_scale/<run>/eval_<step>/*/eval/eval_metrics.json
% (eval/fid_cfg_0.0_sde, eval/fid/clipscore_t2i_cfg_0.0_sde,
%  eval/text_decode/genppl_cfg_0.0_sde, eval/text_decode/clipscore_i2t_cfg_0.0_sde).
\paragraph{Shared versus separate models at small scale.}\label{app:shared-vs-separate-small}
We also compare joint and separate training with a smaller backbone, DiTXA-B/2 (12 blocks,
hidden size 768, 201M parameters). The T2I specialist, the I2T specialist, and the shared
model all use the Table~1 recipe, except that image REPA is applied at layer 4. That keeps
REPA at one third of the trunk's depth, as layer 8 of 24 does in Table~1. All three train
for 200K steps with batch 512 on the same data order. Each step of the shared model sums
the two directional losses on the same batch, so at every checkpoint the shared model has
used the same number of examples and optimizer updates per direction as each specialist,
and the same total compute as the two specialists combined. It does this with half their
combined parameters. We evaluate at 100K and 200K steps with $\omega=0$, 500 SDE steps, and
10K held-out samples per direction. These FIDs use fewer samples than Table~1, so they are
comparable only within this table.

\begin{table}[t]
  \centering
  \small
  \begin{tabular}{llcccc}
    \toprule
    Model (DiTXA-B/2) & Steps & FID$_{10\mathrm{K}}$ ($\downarrow$) & CLIP T2I ($\uparrow$) & Gen PPL ($\downarrow$) & CLIP I2T ($\uparrow$) \\
    \midrule
    Specialists (T2I / I2T) & 100K & \tbd & \tbd & \tbd & \tbd \\
    Shared & 100K & \tbd & \tbd & \tbd & \tbd \\
    \midrule
    Specialists (T2I / I2T) & 200K & \tbd & \tbd & \tbd & \tbd \\
    Shared & 200K & \tbd & \tbd & \tbd & \tbd \\
    \bottomrule
  \end{tabular}
  \caption{\textbf{One bidirectional network versus two specialists at small scale}
    (compute- and data-matched per direction; $\omega=0$).}
  \label{tab:shared-vs-separate-small}
\end{table}
